%%%%%%%%%%%%%%%%
%%% exod 16 %%%
%%%%%%%%%%%%%%%%

\Chapter{16}

\Verse{1}
{
	\hP{וַיִּסְעוּ מֵאֵילִם וַיָּבֹאוּ כּל עֲדַת בְּנֵי יִשְׂרָאֵל אֶל מִדְבַּר סִין אֲשֶׁר בֵּין אֵילִם וּבֵין סִינָי בַּחֲמִשָּׁה עָשָׂר יוֹם לַחֹדֶשׁ הַשֵּׁנִי לְצֵאתָם מֵאֶרֶץ מִצְרָיִם׃}
}
{
	\eP{
		16 And they traveled from Elim, and all of the\\
		congregation of the children of Israel came to the\\
		wilderness of Sîn, which is between Elim and Sinai, on the\\
		fifteenth day of the second month after their exodus from\\
		the land of Egypt.\\
	}
}
{
}

\Verse{2}
{
	\hP{(וילינו) [וַיִּלּוֹנוּ] כּל עֲדַת בְּנֵי יִשְׂרָאֵל עַל מֹשֶׁה וְעַל אַהֲרֹן בַּמִּדְבָּר׃}
}
{
	\eP{
		And all of the congregation of the children of Israel\\
		complained at Moses and at Aaron in the wilderness.\\
	}
}
{
%	\footnote{
%		at Moses and at Aaron. This is one of the most explicit cases in the
%		Bible of the way humans focus on the visible human figure rather than on
%		God. They complain at Moses and Aaron. They say, “You brought us out”
%		(16:3). Moses answers that they will soon see that “YHWH brought you
%		out” (16:6). And he tells them that their complaints are “at YHWH,” and
%		he adds, “What are we, that you complain at us?!” (16:8). The miracle of
%		manna follows, confirming the message that Moses has just told them in
%		three different ways. But by the very next chapter: “the people
%		complained at Moses and said, ‘Why did you bring us up?’” (17:3). They
%		continue to resist the fact that they are experiencing divine power,
%		that they are actually encountering God.
%	}
}

\Verse{3}
{
	\hP{וַיֹּאמְרוּ אֲלֵהֶם בְּנֵי יִשְׂרָאֵל מִי יִתֵּן מוּתֵנוּ בְיַד יְהֹוָה בְּאֶרֶץ מִצְרַיִם בְּשִׁבְתֵּנוּ עַל סִיר הַבָּשָׂר בְּאכְלֵנוּ לֶחֶם לָשֹׂבַע כִּי הוֹצֵאתֶם אֹתָנוּ אֶל הַמִּדְבָּר הַזֶּה לְהָמִית אֶת כּל הַקָּהָל הַזֶּה בָּרָעָב׃}
}
{
	\eP{
		And the children of Israel said to them, “Who would make\\
		it so that we had died by YHWH’s hand in the land of\\
		Egypt, when we sat by a pot of meat, when we ate bread to\\
		the full! Because you brought us out to this wilderness to\\
		kill this whole community with starvation!”\\
	}
}
{
}

\Verse{4}
{
	\hP{וַיֹּאמֶר יְהֹוָה אֶל מֹשֶׁה הִנְנִי מַמְטִיר לָכֶם לֶחֶם מִן הַשָּׁמָיִם וְיָצָא הָעָם וְלָקְטוּ דְּבַר יוֹם בְּיוֹמוֹ לְמַעַן אֲנַסֶּנּוּ הֲיֵלֵךְ בְּתוֹרָתִי אִם לֹא׃}
}
{
	\eP{
		And YHWH said to Moses, “Here, I’m raining bread from the\\
		skies for you, and let the people go out and gather the\\
		daily ration in its day in order that I can test them:\\
		will they go by my instruction or not.\\
	}
}
{
}

\Verse{5}
{
	\hP{וְהָיָה בַּיּוֹם הַשִּׁשִּׁי וְהֵכִינוּ אֵת אֲשֶׁר יָבִיאוּ וְהָיָה מִשְׁנֶה עַל אֲשֶׁר יִלְקְטוּ יוֹם יוֹם׃}
}
{
	\eP{
		And it will be that on the sixth day they’ll prepare what\\
		they bring in, and it will be twice as much as they’ll\\
		gather on regular days.”\\
	}
}
{
%	\footnote{
%		on regular days. Hebrew yôm yôm; meaning that the sixth day’s yield will
%		be double what is gathered on an ordinary day of the week (Sunday
%		through Thursday).
%	}
}

\Verse{6}
{
	\hP{וַיֹּאמֶר מֹשֶׁה וְאַהֲרֹן אֶל כּל בְּנֵי יִשְׂרָאֵל עֶרֶב וִידַעְתֶּם כִּי יְהֹוָה הוֹצִיא אֶתְכֶם מֵאֶרֶץ מִצְרָיִם׃}
}
{
	\eP{
		And Moses and Aaron said to all the children of Israel,\\
		“At evening you’ll know that YHWH brought you out from the\\
		land of Egypt,\\
	}
}
{
}

\Verse{7}
{
	\hP{וּבֹקֶר וּרְאִיתֶם אֶת כְּבוֹד יְהֹוָה בְּשׁמְעוֹ אֶת תְּלֻנֹּתֵיכֶם עַל יְהֹוָה וְנַחְנוּ מָה כִּי (תלונו) [תַלִּינוּ] עָלֵינוּ׃}
}
{
	\eP{
		and in the morning you’ll see YHWH’s glory, because He has\\
		heard your complaints about YHWH. And what are we, that\\
		you complain at us?!”\\
	}
}
{
%	\footnote{
%		YHWH’s glory. The glory surrounds God when in the presence of humans. It
%		is some sort of indescribable substance that veils or masks the deity,
%		who cannot be seen. “A human will not see me and live.” Even the glory
%		itself is seen through a cloud. This is the first mention of the divine
%		glory, and a few verses later it is seen for the first time. It then
%		does not occur again until the revelation at Sinai. It is impressive,
%		but also sad, that its first appearance comes in response to the
%		people’s complaints over food rather than as the surrounding of the
%		Sinai event. It suggests the extreme seriousness of what is taking place
%		in this first period following the exodus. The people are unprepared,
%		vulnerable, and frightened; and they are apparently in need of manifest
%		evidence that the power that brought them out of Egypt is still present,
%		that the exodus was just the beginning, that what comes next is the
%		acquaintance and continuing relationship with God.
%	}
}

\Verse{8}
{
	\hP{וַיֹּאמֶר מֹשֶׁה בְּתֵת יְהֹוָה לָכֶם בָּעֶרֶב בָּשָׂר לֶאֱכֹל וְלֶחֶם בַּבֹּקֶר לִשְׂבֹּעַ בִּשְׁמֹעַ יְהֹוָה אֶת תְּלֻנֹּתֵיכֶם אֲשֶׁר אַתֶּם מַלִּינִם עָלָיו וְנַחְנוּ מָה לֹא עָלֵינוּ תְלֻנֹּתֵיכֶם כִּי עַל יְהֹוָה׃}
}
{
	\eP{
		And Moses said, “When YHWH gives you meat to eat in the\\
		evening and bread to the full in the morning, because YHWH\\
		has heard your complaints that you’re making at Him! And\\
		what are we? Your complaints are not at us but at YHWH!”\\
	}
}
{
}

\Verse{9}
{
	\hP{וַיֹּאמֶר מֹשֶׁה אֶל אַהֲרֹן אֱמֹר אֶל כּל עֲדַת בְּנֵי יִשְׂרָאֵל קִרְבוּ לִפְנֵי יְהֹוָה כִּי שָׁמַע אֵת תְּלֻנֹּתֵיכֶם׃}
}
{
	\eP{
		And Moses said to Aaron, “Say to all of the congregation\\
		of the children of Israel, ‘Come close in front of YHWH,\\
		because He has heard your complaints.’”\\
	}
}
{
}

\Verse{10}
{
	\hP{וַיְהִי כְּדַבֵּר אַהֲרֹן אֶל כּל עֲדַת בְּנֵי יִשְׂרָאֵל וַיִּפְנוּ אֶל הַמִּדְבָּר וְהִנֵּה כְּבוֹד יְהֹוָה נִרְאָה בֶּעָנָן׃}
}
{
	\eP{
		And it was as Aaron was speaking to all of the\\
		congregation of the children of Israel: and they turned to\\
		the wilderness, and here was YHWH’s glory appearing in a\\
		cloud.\\
	}
}
{
}

\Verse{11}
{
	\hP{וַיְדַבֵּר יְהֹוָה אֶל מֹשֶׁה לֵּאמֹר׃}
}
{
	\eP{And YHWH spoke to Moses, saying,}
}
{
}

\Verse{12}
{
	\hP{שָׁמַעְתִּי אֶת תְּלוּנֹּת בְּנֵי יִשְׂרָאֵל דַּבֵּר אֲלֵהֶם לֵאמֹר בֵּין הָעַרְבַּיִם תֹּאכְלוּ בָשָׂר וּבַבֹּקֶר תִּשְׂבְּעוּ לָחֶם וִידַעְתֶּם כִּי אֲנִי יְהֹוָה אֱלֹהֵיכֶם׃}
}
{
	\eP{
		“I’ve heard the complaints of the children of Israel.\\
		Speak to them, saying, ‘Between the two evenings you shall\\
		eat meat, and in the morning you shall be filled with\\
		bread, and you will know that I am YHWH, your God.’”\\
	}
}
{
}

\Verse{13}
{
	\hP{וַיְהִי בָעֶרֶב וַתַּעַל הַשְּׂלָו וַתְּכַס אֶת הַמַּחֲנֶה וּבַבֹּקֶר הָיְתָה שִׁכְבַת הַטַּל סָבִיב לַמַּחֲנֶה׃}
}
{
	\eP{
		And it was in the evening, and quail went up and covered\\
		the camp. And in the morning there was a layer of dew\\
		around the camp.\\
	}
}
{
}

\Verse{14}
{
	\hP{וַתַּעַל שִׁכְבַת הַטָּל וְהִנֵּה עַל פְּנֵי הַמִּדְבָּר דַּק מְחֻסְפָּס דַּק כַּכְּפֹר עַל הָאָרֶץ׃}
}
{
	\eP{
		And the layer of dew went up, and here, on the face of the\\
		wilderness: scaly thin, thin as frost on the ground.\\
	}
}
{
}

\Verse{15}
{
	\hP{וַיִּרְאוּ בְנֵי יִשְׂרָאֵל וַיֹּאמְרוּ אִישׁ אֶל אָחִיו מָן הוּא כִּי לֹא יָדְעוּ מַה הוּא וַיֹּאמֶר מֹשֶׁה אֲלֵהֶם הוּא הַלֶּחֶם אֲשֶׁר נָתַן יְהֹוָה לָכֶם לְאכְלָה׃}
}
{
	\eP{
		And the children of Israel saw, and each said to his\\
		brother, “What is it?”—because they did not know what it\\
		was. And Moses said to them, “That is the bread that YHWH\\
		has given you for food.\\
	}
}
{
%	\footnote{
%		What is it? Hebrew mn hû!. This is a play on words, because these words
%		can also mean “It is manna.” It thus explains the origin of the word
%		“manna” as coming from their not having a name for something new. The
%		famous comedy routine of “Who’s on first” in English is actually an
%		equivalent that conveys what is happening here. They say, “It’s what?”
%		And people then take this not as a question but as a statement: “It’s
%		‘What’!” And it is known thereafter as “What.”
%	}
}

\Verse{16}
{
	\hP{זֶה הַדָּבָר אֲשֶׁר צִוָּה יְהֹוָה לִקְטוּ מִמֶּנּוּ אִישׁ לְפִי אכְלוֹ עֹמֶר לַגֻּלְגֹּלֶת מִסְפַּר נַפְשֹׁתֵיכֶם אִישׁ לַאֲשֶׁר בְּאהֳלוֹ תִּקָּחוּ׃}
}
{
	\eP{
		This is the thing that YHWH has commanded: Collect some of\\
		it, each according to what he eats; you shall take an omer\\
		per head, the number of your persons, each for whoever is\\
		in his tent.”\\
	}
}
{
}

\Verse{17}
{
	\hP{וַיַּעֲשׂוּ כֵן בְּנֵי יִשְׂרָאֵל וַיִּלְקְטוּ הַמַּרְבֶּה וְהַמַּמְעִיט׃}
}
{
	\eP{
		And the children of Israel did so. And they collected—the\\
		one who took the most and the one who took the least.\\
	}
}
{
}

\Verse{18}
{
	\hP{וַיָּמֹדּוּ בָעֹמֶר וְלֹא הֶעְדִּיף הַמַּרְבֶּה וְהַמַּמְעִיט לֹא הֶחְסִיר אִישׁ לְפִי אכְלוֹ לָקָטוּ׃}
}
{
	\eP{
		And they measured by the omer, and the one who took the\\
		most had not exceeded, and the one who took the least had\\
		not fallen short. Each had collected according to what he\\
		eats.\\
	}
}
{
}

\Verse{19}
{
	\hP{וַיֹּאמֶר מֹשֶׁה אֲלֵהֶם אִישׁ אַל יוֹתֵר מִמֶּנּוּ עַד בֹּקֶר׃}
}
{
	\eP{
		And Moses said to them, “Let no one leave any of it over\\
		until morning.”\\
	}
}
{
}

\Verse{20}
{
	\hP{וְלֹא שָׁמְעוּ אֶל מֹשֶׁה וַיּוֹתִרוּ אֲנָשִׁים מִמֶּנּוּ עַד בֹּקֶר וַיָּרֻם תּוֹלָעִים וַיִּבְאַשׁ וַיִּקְצֹף עֲלֵהֶם מֹשֶׁה׃}
}
{
	\eP{
		And they did not listen to Moses, and people left some of\\
		it over until morning, and it yielded worms and stank. And\\
		Moses was angry at them.\\
	}
}
{
}

\Verse{21}
{
	\hP{וַיִּלְקְטוּ אֹתוֹ בַּבֹּקֶר בַּבֹּקֶר אִישׁ כְּפִי אכְלוֹ וְחַם הַשֶּׁמֶשׁ וְנָמָס׃}
}
{
	\eP{
		And they collected it morning by morning, each by what he\\
		eats, and when the sun was hot it melted.\\
	}
}
{
}

\Verse{22}
{
	\hP{וַיְהִי בַּיּוֹם הַשִּׁשִּׁי לָקְטוּ לֶחֶם מִשְׁנֶה שְׁנֵי הָעֹמֶר לָאֶחָד וַיָּבֹאוּ כּל נְשִׂיאֵי הָעֵדָה וַיַּגִּידוּ לְמֹשֶׁה׃}
}
{
	\eP{
		And it was: on the sixth day they collected double bread,\\
		two times the omer for each one, and all the chiefs of the\\
		congregation came and told Moses.\\
	}
}
{
}

\Verse{23}
{
	\hP{וַיֹּאמֶר אֲלֵהֶם הוּא אֲשֶׁר דִּבֶּר יְהֹוָה שַׁבָּתוֹן שַׁבַּת קֹדֶשׁ לַיהֹוָה מָחָר אֵת אֲשֶׁר תֹּאפוּ אֵפוּ וְאֵת אֲשֶׁר תְּבַשְּׁלוּ בַּשֵּׁלוּ וְאֵת כּל הָעֹדֵף הַנִּיחוּ לָכֶם לְמִשְׁמֶרֶת עַד הַבֹּקֶר׃}
}
{
	\eP{
		And he said to them, “That is what YHWH spoke. Tomorrow is\\
		a ceasing, a holy Sabbath to YHWH. Bake what you’ll bake,\\
		and cook what you’ll cook, and leave what is left over for\\
		keeping until the morning.”\\
	}
}
{
}

\Verse{24}
{
	\hP{וַיַּנִּיחוּ אֹתוֹ עַד הַבֹּקֶר כַּאֲשֶׁר צִוָּה מֹשֶׁה וְלֹא הִבְאִישׁ וְרִמָּה לֹא הָיְתָה בּוֹ׃}
}
{
	\eP{
		And they left it until the morning, as Moses had\\
		commanded, and it did not stink, and there was not a worm\\
		in it.\\
	}
}
{
%	\footnote{
%		as Moses had commanded. Until now, the text has always said “as YHWH had
%		commanded” or “as God had commanded.” Here, for the first time, it
%		attributes the commanding to Moses—even though Moses has said explicitly
%		in the preceding verse that this is what YHWH has spoken. Moses still
%		acts only by God’s direction, but the wording participates in the
%		gradual impression through the Tanak that humans’ share in the
%		responsibility for their fate is growing.
%	}
}

\Verse{25}
{
	\hP{וַיֹּאמֶר מֹשֶׁה אִכְלֻהוּ הַיּוֹם כִּי שַׁבָּת הַיּוֹם לַיהֹוָה הַיּוֹם לֹא תִמְצָאֻהוּ בַּשָּׂדֶה׃}
}
{
	\eP{
		And Moses said, “Eat it today, because today is a Sabbath\\
		to YHWH. Today you won’t find it in the field.\\
	}
}
{
}

\Verse{26}
{
	\hP{שֵׁשֶׁת יָמִים תִּלְקְטֻהוּ וּבַיּוֹם הַשְּׁבִיעִי שַׁבָּת לֹא יִהְיֶה בּוֹ׃}
}
{
	\eP{
		Six days you shall collect it. And on the seventh day is a\\
		Sabbath; there will not be any in it.”\\
	}
}
{
}

\Verse{27}
{
	\hP{וַיְהִי בַּיּוֹם הַשְּׁבִיעִי יָצְאוּ מִן הָעָם לִלְקֹט וְלֹא מָצָאוּ׃}
}
{
	\eP{
		And it was: on the seventh day some of the people went out\\
		to collect, and they did not find any.\\
	}
}
{
}

\Verse{28}
{
	\hP{וַיֹּאמֶר יְהֹוָה אֶל מֹשֶׁה עַד אָנָה מֵאַנְתֶּם לִשְׁמֹר מִצְוֺתַי וְתוֹרֹתָי׃}
}
{
	\eP{
		And YHWH said to Moses, “How long do you refuse to observe\\
		my commandments and my instructions?\\
	}
}
{
}

\Verse{29}
{
	\hP{רְאוּ כִּי יְהֹוָה נָתַן לָכֶם הַשַּׁבָּת עַל כֵּן הוּא נֹתֵן לָכֶם בַּיּוֹם הַשִּׁשִּׁי לֶחֶם יוֹמָיִם שְׁבוּ אִישׁ תַּחְתָּיו אַל יֵצֵא אִישׁ מִמְּקֹמוֹ בַּיּוֹם הַשְּׁבִיעִי׃}
}
{
	\eP{
		See that YHWH has given you the Sabbath. On account of\\
		this He is giving you two days’ bread on the sixth day.\\
		Stay, each in his place. Let no man go out from his place\\
		in the seventh day.”\\
	}
}
{
}

\Verse{30}
{
	\hP{וַיִּשְׁבְּתוּ הָעָם בַּיּוֹם הַשְּׁבִעִי׃}
}
{
	\eP{And the people ceased in the seventh day.}
}
{
}

\Verse{31}
{
	\hP{וַיִּקְרְאוּ בֵית יִשְׂרָאֵל אֶת שְׁמוֹ מָן וְהוּא כְּזֶרַע גַּד לָבָן וְטַעְמוֹ כְּצַפִּיחִת בִּדְבָשׁ׃}
}
{
	\eP{
		And the house of Israel called its name “manna.” And it\\
		was like a coriander seed, white, and its taste was like a\\
		wafer in honey.\\
	}
}
{
}

\Verse{32}
{
	\hP{וַיֹּאמֶר מֹשֶׁה זֶה הַדָּבָר אֲשֶׁר צִוָּה יְהֹוָה מְלֹא הָעֹמֶר מִמֶּנּוּ לְמִשְׁמֶרֶת לְדֹרֹתֵיכֶם לְמַעַן יִרְאוּ אֶת הַלֶּחֶם אֲשֶׁר הֶאֱכַלְתִּי אֶתְכֶם בַּמִּדְבָּר בְּהוֹצִיאִי אֶתְכֶם מֵאֶרֶץ מִצְרָיִם׃}
}
{
	\eP{
		And Moses said, “This is the thing that YHWH has\\
		commanded: ‘an omer-ful of it is for watching over through\\
		your generations so that they will see the bread that I\\
		fed you in the wilderness when I brought you out from the\\
		land of Egypt.’”\\
	}
}
{
}

\Verse{33}
{
	\hP{וַיֹּאמֶר מֹשֶׁה אֶל אַהֲרֹן קַח צִנְצֶנֶת אַחַת וְתֶן שָׁמָּה מְלֹא הָעֹמֶר מָן וְהַנַּח אֹתוֹ לִפְנֵי יְהֹוָה לְמִשְׁמֶרֶת לְדֹרֹתֵיכֶם׃}
}
{
	\eP{
		And Moses said to Aaron, “Take one jar and place an\\
		omer-ful of manna there and lay it in front of YHWH for\\
		watching over through your generations,”\\
	}
}
{
}

\Verse{34}
{
	\hP{כַּאֲשֶׁר צִוָּה יְהֹוָה אֶל מֹשֶׁה וַיַּנִּיחֵהוּ אַהֲרֹן לִפְנֵי הָעֵדֻת לְמִשְׁמָרֶת׃}
}
{
	\eP{
		as YHWH had commanded Moses. And Aaron laid it in front of\\
		the Testimony for watching over.\\
	}
}
{
%	\footnote{
%		the Testimony. This term is introduced later as referring to the words
%		of the Ten Commandments on the “tablets of the Testimony” (Exod 31:18),
%		which are kept in the “Ark of the Testimony” (25:22), which is in turn
%		housed in the “Tabernacle of the Testimony” (38:21). It seems
%		chronologically out of place to report here that Aaron placed the manna
%		there, when the tablets, ark, and Tabernacle have not yet been made.
%		This problem has nothing to do with questions of authorship in critical
%		biblical scholarship. Nor is it to be explained by the traditional
%		rabbinic principle of “There is no early or late in the Torah.” Rather,
%		there is a series of items that are reported to have been placed by the
%		Testimony in the Tabernacle: the manna, the incense, Aaron’s staff. The
%		placement of each is reported in the context of the story or law that
%		relates to it. Since the story of the manna is told here, the report
%		that Aaron placed it at the Testimony is included here as well. To have
%		saved this report for a later, more chronologically appropriate, point
%		in the Torah would have meant cutting it off from the rest of the manna
%		story. The report in the next verse that they ate the manna for forty
%		years likewise goes chronologically past the present moment in the
%		story. It thus further confirms that the purpose of the present context
%		is to cover the entire matter of the manna in one place.
%	}
}

\Verse{35}
{
	\hP{וּבְנֵי יִשְׂרָאֵל אָכְלוּ אֶת הַמָּן אַרְבָּעִים שָׁנָה עַד בֹּאָם אֶל אֶרֶץ נוֹשָׁבֶת אֶת הַמָּן אָכְלוּ עַד בֹּאָם אֶל קְצֵה אֶרֶץ כְּנָעַן׃}
}
{
	\eP{
		And the children of Israel ate the manna forty years until\\
		they came to settled land. They ate the manna until they\\
		came to the edge of the land of Canaan.\\
	}
}
{
}

\Verse{36}
{
	\hP{וְהָעֹמֶר עֲשִׂרִית הָאֵיפָה הוּא׃}
}
{
	\eP{(And an omer is a tenth of an ephah.)}
}
{
}
